%% notes-tex/010-index-defect/sections/2-evidence.tex

\section{The evidence\secstatus{todo}}
\label{sec:evidence}

Both defects are confirmed rather than inferred. The confirmation matters because
the first symptom, one checkpoint appearing to disagree with itself across runs,
has an obvious wrong reading: that the model is nondeterministic. It is not.

\subsection{The symptom\secstatus{todo}}
\label{sec:symptom}

Checkpoint \texttt{c7671wgj} scored three separate design spaces:
\texttt{inference\_1}, 4{,}370{,}595 triples over a 526-gene panel;
\texttt{inference\_2}, 479{,}195 triples; and \texttt{inference\_3},
465{,}735{,}532 triples genome-wide, the space the panel-12 and panel-24
selections consumed. A triple scored in two of them should receive the same
number twice, and a gene's mean predicted effect, averaged over hundreds of
triples, should agree between runs.

%% SOURCE: experiments/010-kuzmin-tmi/scripts/inference_run_consistency.py,
%% results/inference_run_consistency.csv
\begin{table}[htbp]
\centering
\begin{tabular}{llr}
\hline
Comparison & $n$ & Pearson \\
\hline
Same triple, \texttt{inference\_2} vs \texttt{inference\_3} & 2{,}404 triples & 1.0000 \\
Same triple, \texttt{inference\_1} vs \texttt{inference\_3} & 330 triples & 0.295 \\
Per-gene mean, \texttt{inference\_2} vs \texttt{inference\_3} & 1{,}274 genes & 0.800 \\
Per-gene mean, \texttt{inference\_1} vs \texttt{inference\_3} & 395 genes & 0.037 \\
Per-gene mean, \texttt{inference\_1} vs \texttt{inference\_2} & 179 genes & 0.167 \\
\hline
\end{tabular}
\caption{Agreement of checkpoint \texttt{c7671wgj} with itself across inference
runs. Runs 2 and 3 are identical on shared triples, mean absolute difference
$0.000000$, so they share a pipeline. Run 1 is unrelated to them. Per-gene means
average over hundreds of triples each, so sampling noise cannot explain
$0.037$.}
\label{tab:runconsistency}
\end{table}

Deciding which run handles gene identity correctly needs a reference outside the
transformer. The per-gene coefficient $\hat{\beta}_g$ of the additive ridge fit
in the companion document (\texttt{notes-tex/010-additive-baselines}) is one: it
is estimated from measured training labels, keyed on systematic gene name, and
independently validated at 0.400 test Pearson. If a run's predictions carry
real gene identity, a gene's mean predicted effect should track
$\hat{\beta}_g$.

%% SOURCE: experiments/010-kuzmin-tmi/scripts/inference_run_consistency.py
\begin{table}[htbp]
\centering
\begin{tabular}{lrr}
\hline
Run & Own gene set & 167 common genes \\
\hline
\texttt{inference\_1} & 0.579 (383 genes) & 0.576 \\
\texttt{inference\_2} & 0.048 (1{,}138 genes) & 0.231 \\
\texttt{inference\_3} & $-0.008$ (3{,}180 genes) & 0.174 \\
\hline
\end{tabular}
\caption{Pearson correlation between a run's per-gene mean prediction and the
additive ridge coefficient for that gene. The right column restricts to the 167
genes all three runs score with at least 200 observations, so the different gene
coverage of the three runs cannot explain the gap.}
\label{tab:anchor}
\end{table}

Restricted to identical genes, \texttt{inference\_1} tracks gene identity more
than three times as strongly as the two runs that produced the selection. The
same comparison made across a whole design space says the same thing: over the
325{,}631{,}845 \texttt{inference\_3} triples whose three genes all carry an
additive coefficient, the additive model and the checkpoint correlate at
$0.0018$ with a top-100 overlap of zero, while on \texttt{inference\_1} the three
checkpoints correlate with it at $0.548$, $0.579$ and $0.511$.

\tcfigfit[0.92]{figures/inference_run_consistency.pdf}{Per-gene mean prediction
against the additive ridge coefficient, for each of the three inference runs of
checkpoint \texttt{c7671wgj}. The additive coefficient is fit on measured
training data and keyed on systematic gene name, so it anchors gene identity
from outside the transformer.}{fig:anchor}

\subsection{The reference the test is anchored on\secstatus{todo}}

A hand-built forward pass that computes the encoder once and feeds it only
perturbed-gene indices reproduces the recorded validation metrics of all three
checkpoints to the fourth decimal, for example $0.461917$ against a recorded
$0.461881$. That agreement is what licenses the path as a reference: a wrong gene
ordering or a mis-loaded weight would give a plausible but wrong number instead.

\subsection{The decisive test\secstatus{todo}}

Two independent 384-triple samples were taken from each stored prediction file,
at different offsets, and re-scored through that reference under each candidate
index map.

%% SOURCE: results/offby28_index_verification.csv, written by the verification
%% run described in experiments.010-kuzmin-tmi.scripts.rescore_panel_triples_corrected
\begin{table}[htbp]
\centering
\begin{tabular}{llrr}
\hline
Run & sample & 6{,}607-gene map & 6{,}579-gene map \\
\hline
\texttt{inference\_1} & head & \textbf{0.999998} & 0.115 \\
\texttt{inference\_1} & middle & \textbf{0.999999} & $-0.252$ \\
\texttt{inference\_3} & head & 0.074 & \textbf{0.999999} \\
\texttt{inference\_3} & middle & $-0.003$ & \textbf{0.999998} \\
\hline
\end{tabular}
\caption{Pearson correlation between each run's stored predictions and the
validated scoring path, under each candidate index space. Every run is
reproducible under exactly one map and neither is reproducible under both. Mean
absolute difference at the reproducing map is about $2 \times 10^{-5}$, which is
the runs' half-precision autocast.}
\label{tab:indextest}
\end{table}

Three things follow. The model is deterministic, since every stored prediction is
exactly recoverable. \texttt{inference\_1} is correct and
\texttt{inference\_2} and \texttt{inference\_3} are not. And the perfect agreement
between runs 2 and 3 that first looked reassuring was evidence they shared the
same wrong map, not evidence either was right.

\subsection{What is invalid and what is not\secstatus{todo}}

Invalid: everything derived from the \texttt{inference\_2} and
\texttt{inference\_3} prediction files, which is the panel-12 and panel-24 gene
selections and every table and figure built from them.

Unaffected: \texttt{inference\_1}, whose 526-gene panel lies entirely inside the
model's gene set, so it escaped both defects rather than only the first. Also
unaffected is every number computed on the labeled split through the validated
path, including the held-out metrics and the additive-null comparison.

\subsection{The guard, and the fix it is not\secstatus{todo}}
\label{sec:guard}

The model now raises when the cell graph's gene node count differs from the
\texttt{gene\_num} the checkpoint was built for. That converts a silent
relabeling of the entire genome into a failure at the first forward pass.

It is detection, not prevention. The condition it catches is a symptom of the
design in Section~\ref{sec:rebuilt}, where the index space is recomputed from a
live genome at every use. Pinning the ordering with the build and loading it
would make the failure impossible instead of merely loud. A commit hash cannot
substitute: the library was untouched between the two runs, and what moved was
the genome underneath them, so two runs at the same commit scored different
genes.
